\documentclass[11pt]{article}

\usepackage{times}
\usepackage{float}
\renewcommand*{\ttdefault}{cmtt}
\usepackage{fullpage}
\usepackage[compact]{titlesec}  
\titlespacing{\section}{0pt}{0pt}{0pt}
\titleformat{\section}
  {\normalfont\fontsize{13}{15}\bfseries}{\thesection}{1em}{}
\titleformat{\section}
  {\normalfont\fontsize{13}{15}\bfseries}{\thesection}{1em}{}
\usepackage[hidelinks]{hyperref}

\usepackage[utf8]{inputenc}
\usepackage{booktabs}
\usepackage{subcaption}
\usepackage{caption}
\usepackage{enumitem}
\usepackage{authblk}
\usepackage{caption}
\captionsetup[table]{skip=3pt}
\usepackage{multirow}
\usepackage{wrapfig}
\usepackage[round]{natbib}
\usepackage{textcomp}
\usepackage{adjustbox}
\usepackage{tabularx}
\usepackage{todonotes}
\usepackage{footnote}
\usepackage{graphicx}
\usepackage{fancyvrb} 
\usepackage{hyperref}
\hypersetup{linkcolor=blue}
\usepackage{listings}
\usepackage{placeins}
\usepackage{amsmath}
\usepackage{amssymb}
\usepackage{pifont}
\usepackage{comment}

\lstset{
  language=Python,
  aboveskip=3mm,
  belowskip=3mm,
  showstringspaces=false,
  columns=flexible,
  basicstyle={\small\ttfamily},
  numbers=none,
  numberstyle=\tiny\color{gray},
  keywordstyle=\color{blue},
  commentstyle=\color{dkgreen},
  stringstyle=\color{mauve},
  breaklines=true,
  breakatwhitespace=true,
  tabsize=4
}

\newcommand{\shrimai}[1]{{\color{red}{\textbf{Shrimai}: #1}}}
\newcommand{\jupinder}[1]{{\color{blue}{\textbf{Jupinder}: #1}}}
\newcommand{\joseph}[1]{{\color{red}{\textbf{Joseph}: #1}}}

\newcommand{\ours}{Nemotron-4}

\title{Nemotron-4 15B Technical Report}
\author{Jupinder Parmar\footnote{
Equal contribution, corresponding authors: {\tt 
\{jupinderp,sprabhumoye,jjennings,mpatwary\}@nvidia.com}. 
} \ \ Shrimai Prabhumoye$^{*}$ \ \ Joseph Jennings$^{*}$ \ \ Mostofa Patwary$^{*}$ \\ Sandeep Subramanian\footnote{Work done while at NVIDIA.} \ \ Dan Su \ \ Chen Zhu \ \ Deepak Narayanan \ \  Aastha Jhunjhunwala \ \  Ayush Dattagupta \ \ Vibhu Jawa \ \ Jiwei Liu \ \ Ameya Mahabaleshwarkar \ \ Osvald Nitski \ \  Annika Brundyn \ \ James Maki \ \ Miguel Martinez \ \ Jiaxuan You \ \ John Kamalu \ \ Patrick LeGresley \ \ Denys Fridman \ \ Jared Casper \ \  Ashwath Aithal \ \ Oleksii Kuchaiev \ \ Mohammad Shoeybi \ \ Jonathan Cohen \ \ Bryan Catanzaro \par \textbf{NVIDIA}}

\date{}

\begin{document}

\parindent 0.0in
\parskip 0.15in

\maketitle

\begin{abstract}
We introduce \ours{} 15B, a 15-billion-parameter large multilingual language model trained on 8 trillion text tokens. \ours{} 15B demonstrates strong performance when assessed on English, multilingual, and coding tasks: it outperforms all existing similarly-sized open models on 4 out of 7 downstream evaluation areas and achieves competitive performance to the leading open models in the remaining ones. Specifically, \ours{} 15B exhibits the best multilingual capabilities of all similarly-sized models, even outperforming models over four times larger and those explicitly specialized for multilingual tasks.

\end{abstract}

\section{Introduction}

Recently published efforts \citep{hoffmann2022training,touvron2023llama,touvron2023llama2,yang2023baichuan, jiang2023mistral} in language model pre-training have been inspired by Chinchilla scaling laws \citep{hoffmann2022training}, which argue for scaling data along with model size given a fixed compute budget, compared to past work that only scaled the size of the model~\citep{kaplan2020scaling,brown2020language, smith2022,rae2022scaling,workshop2023bloom}.
For example, \citep{hoffmann2022training} shows that given two roughly IsoFLOP GPT models with a similar data distribution, a 65-billion-parameter model on 1.4 trillion tokens and a 280-billion-parameter model on 300 billion tokens, the 65B model has better accuracy on downstream tasks. 

This trade-off of allocating compute towards training on more data as opposed to increasing model size is particularly appealing from an inference perspective, reducing latency and the amount of compute needed to serve models. As a consequence, a major focus of language modeling training efforts has shifted to collecting high-quality multi-trillion token datasets from public sources such as Common Crawl. We continue this trend by introducing \ours{} 15B which was trained on 8 trillion tokens of English, multilingual, and coding text and was developed to be the best general-purpose large language model (LLM) that can fit on a single NVIDIA A100 or H100 GPU.

As demonstrated in Figure \ref{fig:benchmark_comp}, \ours{} 15B exhibits high downstream accuracies across a wide range of English, code, and multilingual evaluation areas. In comparison to leading similarly-sized, open models we show that \ours{} 15B is significantly better than LLaMA-2 34B~\citep{touvron2023llama2}, which has over twice the number of parameters, and is better than Mistral 7B~\citep{jiang2023mistral} on all English evaluation areas. Additionally, \ours{} 15B achieves competitive accuracies to QWEN 14B~\citep{bai2023qwen} and Gemma 7B~\citep{gemma24}. In a comparison across a wide range of programming languages, we find that \ours{} 15B achieves better average accuracy, and in particular on low-resource programming languages, than Starcoder~\citep{li2023starcoder}, a code-specific model, and Mistral 7B. As \ours{} 15B was trained on significant amount of multilingual data, it is currently the state-of-the-art general purpose model in its size class on all multilingual benchmarks. We find that \ours{} is better than PALM 62B-Cont~\citep{palm2}, and also outperforms multilingual-specific models such as XGLM~\citep{lin2022fewshot} and mGPT~\citep{shliazhko2022mgpt}.

\begin{figure}
\centering
\begin{subfigure}{.5\textwidth}
  \centering
  \includegraphics[width=\linewidth]{comparison_largest_fin.pdf}
  \label{fig:sub1}
\end{subfigure}
\begin{subfigure}{.5\textwidth}
  \centering
  \includegraphics[width=\linewidth]{multilingual_comp_largest_fin.pdf}
  \label{fig:sub2}
\end{subfigure}
\caption{Comparison of \ours{} 15B across seven evaluation areas against similarly sized models. The composition of tasks that form each evaluation area can be found, along with more detailed evaluation results, in Section \ref{results}}
\label{fig:benchmark_comp}
\end{figure}

\begin{table}[h]
    \centering
    \begin{tabular}{ccccccc}
    \toprule
    Number of          & Hidden    & Number of       & Number of & Sequence & Vocabulary \\
    transformer layers & dimension & attention heads &    KV heads                & length   & size \\
    \toprule
    32 & 6144 & 48 & 8 & 4096 & 256,000 \\
    \bottomrule
    \end{tabular}
    \caption{Key hyper-parameters affecting size of \ours{} 15B.}
    \label{tab:model_arch}
\end{table}

\section{Architecture Details}

\ours{} uses a standard decoder-only Transformer architecture~\citep{DBLP:journals/corr/VaswaniSPUJGKP17}, with causal attention masks. Exact hyper-parameters affecting size are shown in Table~\ref{tab:model_arch}. 
\ours{} has 3.2 billion embedding parameters and 12.5 billion non-embedding parameters.
We use Rotary Position Embeddings (RoPE)~\citep{su2021roformer}, SentencePiece tokenizer~\citep{kudo2018sentencepiece}, squared ReLU activations in the MLP layers, no bias terms, dropout rate of zero, and untied input-output embeddings.
We use grouped query attention (GQA)~\citep{ainslie2023gqa} for faster inference and lower memory footprint.

\paragraph{Data.} 

We train \ours{} 15B on a pre-training dataset consisting of 8 trillion tokens. At a high-level, the data blend is split into three different types of data: English natural language data (70\%), multilingual natural language data (15\%), and source-code data (15\%). 

\begin{figure}
    \centering
    \includegraphics[width=0.85\linewidth]{plots/new_full_distr.png}
    \caption{Data composition of the English tokens used for pre-training}
    \label{fig:data_comp}
\end{figure}

The English corpus consists of curated documents from a variety of sources and domains including web documents, news articles, scientific papers, books, etc and the distribution used in our pre-training set is highlighted in Figure~\ref{fig:data_comp}. The code and multilingual data consists of a diverse set of natural and programming languages. We find that appropriately sampling tokens from these languages is key to strong accuracies in these domains. We share the distributions used for both code and multilingual tokens in our pre-training dataset in Figure~\ref{fig:code_distr} and Figure~\ref{fig:multi_distr} respectively.  

\begin{figure}[t]
    \centering
    \includegraphics[width=\linewidth]{plots/test_code.png}
    \caption{Data distribution of the 43 programming languages used for pre-training. The number within each bar indicates the percent of the overall code distribution that an individual language comprises.}
    \label{fig:code_distr}
\end{figure}

In constructing the pre-training corpus, we remove any possible duplicates via document-level exact and near-deduplication~\citep{jennings2023ndc}. We additionally applied document-level quality filtering across our corpus using a language-model based filtering approach similar to \citep{wenzek2019ccnet} in addition to a  series of heuristic filters as described in \citep{rae2022scaling} and \citep{raffel2020exploring}.

We train a BPE tokenizer in SentencePiece~\citep{kudo2018sentencepiece} on data that is randomly sampled from the final 8T token dataset. To have better coverage of low-resource languages in the tokenizer, we upsample non-English data relative to the final training dataset distribution. Our tokenizer preserves whitespaces (including leading and trailing ones), splits numbers into their individual digits~\citep{chowdhery2022palm}, and relies on byte-level backoff to handle unknown character sequences. The final vocabulary size is 256,000 tokens. 

\begin{figure}[t]
    \centering
    \includegraphics[width=\linewidth]{plots/test_multi.png}
    \caption{Data distribution of the 53 natural languages, aside from English,we used for pre-training. The number within each bar indicates the percent of the overall multilingual distribution that an individual language comprises.}
    \label{fig:multi_distr}
\end{figure}

\paragraph{Pre-training.} 
\ours{} was trained using 384 DGX H100 nodes; each node contains 8 H100 80GB SXM5 GPUs based on the NVIDIA Hopper architecture~\citep{nvidia2022h100}. 
Each H100 GPU has a peak throughput of 989 teraFLOP/s when doing 16-bit floating point (\texttt{bfloat16}) arithmetic without sparsity. Within each node, GPUs are connected by NVLink and NVSwitch~\citep{nvlink}; the GPU-to-GPU bandwidth is 900 GB/s (450 GB/s in each direction). Each node has 8 NVIDIA Mellanox 400 Gbps HDR InfiniBand Host Channel Adapters (HCAs) for inter-node communication. 

We used a combination of 8-way tensor parallelism~\citep{shoeybi2019megatron} and data parallelism to train the model; we also use a distributed optimizer to shard the optimizer state over the data-parallel replicas. The degree of data parallelism was varied from 96 to 384 as the batch size was ramped up. 
Table~\ref{tab:43bmodel_batch_time_mfu} summarizes the 3 stages of batch size ramp, and includes the per-iteration time and model FLOP/s utilization (MFU)~\citep{chowdhery2022palm,korthikanti2022sequence}. MFU quantifies how efficiently the GPUs are utilized in model training. Training was completed in approximately 13 calendar days.

\begin{table}[H]
    \centering
    \begin{tabular}{crccccc}
    \toprule
    Data-parallel size & GPUs & Iteration time (secs) & MFU (\%) & Batch size & Tokens (B) & Time (days) \\
    \toprule
    96 &  768 & 0.57 & 34.3 & 384 &  200 & 0.8 \\
    192 & 1,536 & 0.58 & 33.3 & 768 &  200 & 0.4 \\
    288 & 2,304 & 0.64 & 30.5 & 1,152 &  7,600 & 11.9   \\
    \bottomrule
    \end{tabular}
    \caption{Batch size rampup schedule, along with time and efficiency metrics for the \ours{} 15B parameter model.}
    \label{tab:43bmodel_batch_time_mfu}
\end{table}

\paragraph{Continued Training.}
Similar to recent work~\citep{geminiteam2023gemini}, we find that switching the data distribution and learning rate decay schedule at the end of model training greatly improves model quality. Concretely, after having trained over the entirety of our 8T pre-training dataset, we use the same loss objective and perform continued training on small number of tokens in comparison to the pre-training tokens. 

In this additional phase of continued training, we utilize two distinct data distributions. 
The first distribution is where the majority of tokens during continued training are sampled from. It utilizes tokens that have already been introduced during pre-training but with a distribution that places larger sampling weight on higher quality sources. 
The second distribution introduces a small number of benchmark-style alignment examples to better allow the model to respond to such questions in downstream evaluations while also up-weighting data sources that come from areas of low model performance. 
In accompaniment with a learning rate schedule that prioritizes a steeper slope of decay than magnitude of learning rate, we find that such an ordering and style of data distributions allows for the model to gently transition from the pre-training dataset and better learn newly emphasized data areas. 

\section{Results}\label{results}
We evaluate \ours{} 15B on a variety of downstream evaluation areas covering a diverse range of tasks and domains. In all evaluations, we adhere to the standardized task setup and share the exact settings used. The covered evaluation categories include: 

\begin{itemize}
    \item \textbf{Commonsense Reasoning (0-shot):}  SIQA \citep{sap2019socialiqa}, ARC easy and challenge \citep{clark2018think},  PIQA \citep{Bisk2020PIQARA}, Winogrande \citep{Sakaguchi2020WINOGRANDEAA}, and Hellaswag  \citep{Zellers2019HellaSwagCA}
    \item \textbf{Popular Aggregated Benchmarks:} MMLU (5-shot) \citep{hendrycks2020measuring} and BBH (3-shot) \citep{suzgun2022challenging}
    \item \textbf{Math:} GSM8K (8-shot with maj@1) \citep{cobbe2021}
    \item \textbf{Code:} Pass@1 scores on HumanEval (0-shot) \citep{chen2021evaluating}, MBPP (3-shot) \citep{austin2021program}, and MultiPL-E (0-shot) \citep{multiple2023}
    \item \textbf{Multilingual:} classification via XCOPA (0 and 4-shot) \citep{xcopa}, machine translation with FLORES-101 (8-shot)~\citep{goyal2021}, and generation tasks such as MGSM (8-shot)~\citep{shi2022language} and TyDiQA (1-shot)~\citep{tydiqa} 
\end{itemize}

In our evaluations, we compare against a number of external decoder-only transformer language models and unless otherwise stated we use the numbers published in the reports of the corresponding models. For English and code tasks, we share detailed results for \ours{} 15B, LlaMA-2 13B and 34B \citep{touvron2023llama2}, Mistral 7B \citep{jiang2023mistral}, Baichuan-2 13B~\citep{yang2023baichuan}, QWEN 14B~\citep{bai2023qwen}, and Gemma 7B \citep{gemma24}. For multilingual benchmarks, we report results against PaLM 62B and 62B-cont \citep{chowdhery2022palm} as well as models specially trained for multilingual capabilities such as mGPT 13B \citep{shliazhko2022mgpt} and XGLM 7.5B \citep{lin2022fewshot}.

\subsection{Commonsense Reasoning}
We use the LM-Evaluation Harness~\citep{eval-harness} to evaluate \ours{} 15B across all aforementioned tasks. Table~\ref{tab:lm_eval_harness_external} showcases that \ours{} 15B achieves the strongest average performance on this diverse set of tasks.

\begin{table}[H]
\centering
\begin{adjustbox}{max width=\textwidth, scale=1.0}
\begin{tabular}{lrccccccc}
\toprule
& Size & SIQA  & ARC-c & ARC-e & PIQA & Winogrande & Hellaswag & AVG \\ 
\midrule
\multirow{2}{*}{LLaMA-2} & 13B & 50.3 & 49.4 & 77.3 & 79.8 & 72.8 & 80.7 & 68.4 \\
& 34B & 50.9 & 54.5 & 79.4 & 81.9 & 76.7 & 83.3 & 71.1 \\
\midrule
\multirow{1}{*}{Baichuan-2} & 13B & - & - & - & 78.1 & - & 70.8 & - \\
\midrule
QWEN & 14B & 77.9 & 84.4 &90.3 & 79.9 & - & 80.2 & - \\
\midrule
Mistral  & 7B & \ \ 47.0$^{*}$ & 55.5 &80.0 & 83.0 & 75.3 & 81.3 & 70.4 \\
\midrule
Gemma & 7B & 51.8 & 53.2 & 81.5 & 81.2 & 72.3 & 81.2 & 70.2 \\
\midrule
\ours{} & 15B & 60.9 & 55.5 & 80.9 & 82.4 & 78.0 & 82.4 & \textbf{73.4} \\
\bottomrule
\end{tabular}
\end{adjustbox}
\caption{\label{tab:lm_eval_harness_external} Results on standard reasoning benchmarks in the zero-shot setting. We report the average across all tasks where possible for a fair comparison. The values marked with $*$ are read from ~\citet{gemma24} }
\end{table}

\begin{comment}
\begin{table}[H]
\centering
\begin{adjustbox}{max width=\textwidth, scale=1.0}
\begin{tabular}{lrccccccccc}
\toprule
& Parameters  & Lambada & ARC-Easy & WebQS & Race-H & ANLI-R2 & PiQA & Winogrande & Hellaswag & AVG \\ 
\midrule
\multirow{2}{*}{LLaMA-2} & 7B & - & 75.2 & - & - & - & 78.8 & 69.2 & 77.2  & 75.1 \\
& 13B & - & 77.3 & - & - & - & 79.8 & 72.8 & 80.7 & 77.7 \\
& 34B & - & 79.4 & & & & 81.9 & 76.7 & 83.3 & 80.3 \\
\midrule
\multirow{2}{*}{Baichuan-2} & 7B & - & - & - & - & - & 76.2 & - & 67.0 & - \\
& 13B & - & - & - & - & - & 78.1 & - & 70.8 & - \\
\midrule
Mistral  & 7B & - & 80.0 & - & - & - & 83.0 & 75.3 & 81.3 & 79.9\\
\midrule
\ours{} & 15B & 75.7 & 80.9 & 22.3 & 47.6 & 44.8 & 82.4 & 78.0 & 82.4 & \textbf{80.9} \\
\bottomrule
\end{tabular}
\end{adjustbox}
\caption{\label{tab:lm_eval_harness_external}Zero-shot accuracy of models on many NLU tasks. We report the average over ARC-Easy,  PiQA, Winogrande, and Hellaswag where possible. }
\end{table}
\end{comment}

\subsection{Popular Aggregated Benchmarks}

The MMLU~\citep{hendrycks2020measuring} and Big Bench Hard (BBH)~\citep{suzgun2022challenging} benchmarks have been developed as a challenging assessment of language models' capabilities on a wide range of tasks and domains. As seen from  Table~\ref{tab:bbh_mmlu}, \ours{} 15B achieves the best score on BBH across existing models at its scale by nearly 7\%. 
Additionally, \ours{} is significantly better than LLaMA-2 70B model on BBH benchmark where LLaMA-2 70B attains a score of 51.2 and \ours{} is 58.7.
\ours{} 15B additionally attains a highly competitive MMLU score and its per-category performance on MMLU can be found in Table \ref{tab:mmlu}.

\begin{table}[h]
\centering
\begin{tabular}{lrcc} 
\toprule
    & Size & BBH & MMLU     \\ 
\midrule
\multirow{3}{*}{LLaMA-2} & 13B  & 39.4 & 54.8 \\
& 34B & 44.1 & 62.6 \\
\midrule
\multirow{1}{*}{Baichuan-2} & 13B  & 48.8 & 59.2 \\
\midrule
QWEN & 14B & 53.4 & \textbf{66.3} \\
\midrule
Mistral & 7B & 39.5 & 60.1 \\
\midrule
Gemma & 7B & 55.1 & 64.3 \\
\midrule
\ours{} & 15B & \textbf{58.7} & 64.2 \\ 
\bottomrule
\end{tabular}
\caption{\label{tab:bbh_mmlu} \ours{} 15B attains highly competitive performance on popular aggregate benchmarks. The BBH result for Mistral is read from the figure in~\citep{jiang2023mistral}.}
\end{table}

\begin{comment}
\begin{table}[H]
\centering
\begin{tabular}{lccccc} 
\toprule
Sub-tasks & \ours{} &  &       \\ 
boolean\_expressions & \\
causal\_judgement  & \\
date\_understanding  & \\
disambiguation\_qa  & \\
dyck\_languages  & \\
formal\_fallacies  & \\
geometric\_shapes  & \\
hyperbaton  & \\
logical\_deduction\_five\_objects  & \\
logical\_deduction\_seven\_objects  & \\
logical\_deduction\_three\_objects  & \\
movie\_recommendation  & \\
multistep\_arithmetic\_two  & \\
navigate  & \\
object\_counting & \\
penguins\_in\_a\_table  & \\
reasoning\_about\_colored\_objects & \\
ruin\_names & \\
salient\_translation\_error\_detection & \\
snarks & \\
sports\_understanding & \\
temporal\_sequences & \\
tracking\_shuffled\_objects\_five\_objects & \\
tracking\_shuffled\_objects\_seven\_objects & \\
tracking\_shuffled\_objects\_three\_objects & \\
web\_of\_lies & \\
word\_sorting & \\
\bottomrule
\end{tabular}
\caption{\label{tab:bbh} Three-shot performance on the Big-Bench Hard.}
\end{table}
\end{comment}

\subsection{Math and Code}

Recently, large language models have been shown to be effective at both mathematical reasoning and a variety of coding tasks \citep{allal2023santacoder,chowdhery2022palm, touvron2023llama}. Table~\ref{tab:code_external} highlights the performance of \ours{} 15B on such tasks. Specifically, on mathematical reasoning we find that \ours{} 15B achieves strong performance as it attains a similar score to Gemma 7B, but lags behind models such as Baichuan-2 and QWEN. On code tasks, we see that \ours{} performs on par with QWEN 14B while remaining slightly behind Gemma 7B. Across both types of tasks, \ours{} 15B is able to outperform Mistral 7B and LlaMA-2 13B/34B. 

\begin{table}[H]
\centering
\begin{adjustbox}{max width=\textwidth, scale=1.0}
\begin{tabular}{lrccc}
\toprule
& Size & GSM8K & HumanEval & MBPP \\
\midrule
  \multirow{2}{*}{LlaMA-2} & 13B & 28.7 & 18.3 &30.6 \\ 
 & 34B & 42.2 & 22.6 & 33.0 \\
 \midrule
 \multirow{1}{*}{Baichuan-2}  & 13B & 52.8 & 17.1  & 30.2 \\ 
 \midrule
 \multirow{1}{*}{QWEN}  & 14B & \textbf{60.1} & 32.2 & 40.8 \\
 \midrule
\multirow{1}{*}{Mistral}  & 7B & 35.4$^{*}$  &  30.5 &  \ \ 40.2$^{*}$   \\
\midrule
\multirow{1}{*}{Gemma}  & 7B & 46.4 &  \textbf{32.3} & \textbf{44.4} \\
\midrule
\multirow{1}{*}{\ours{}}  & 15B & 46.0 &   31.6 & 40.6  \\
\bottomrule
\end{tabular}
\end{adjustbox}
\caption{\label{tab:code_external} Comparative results on math and code benchmarks. As Mistral 7B reports MBPP performance on a different eval split and uses a different evaluation setting for GSM8K , we use the corresponding numbers reported in~\citep{gemma24}}
\end{table}

Nearly all similarly-sized open models determine their code abilities solely based on performance on Python related tasks -- disregarding an evaluation of their capabilities on other programming languages. In Table~\ref{tab:multiple-e}, we demonstrate results of \ours{} 15B on the Multiple-E~\citep{cassano2023multipl} benchmark across 11 diverse programming languages and compare it against Mistral 7B and Starcoder~\citep{li2023starcoder}, a 15B parameter model that has been specially trained for code. We find that \ours{} 15B attains strong coding performance across a wide assortment of programming languages and outperforms both Starcoder and Mistral 7B on average. We especially highlight the superior performance of \ours{} 15B on low-resource programming languages such as Scala, Julia, and R. 

\begin{table}[t]
\centering
\begin{adjustbox}{max width=\textwidth, scale=1.0}
\begin{tabular}{lccccccccccccc}
\toprule
  & Size & JavaScript & Julia & Java & Lua & C++ & C-Sharp & PHP & Shell & TypeScript & R & Scala & AVG  \\
\midrule
Starcoder & 15B & 30.8 & 23.0 & 30.2 & 23.9 & 31.6 & 21.0 & 26.1	& 10.5 & 32.3 & 15.5 & 27.6 & 24.2 \\
Mistral & 7B & 34.2 & 22.0 & 26.0 & 25.3 & 29.1 & 22.8 & 27.9	& 8.9 & 28.5 & 11.8 & 22.2 & 23.6 \\
\ours{} & 15B & 28.6 & 24.8 & 24.8 & 24.2 & 35.4 & 21.1 & 27.3 & 8.9 & 32.9 & 18.6 & 27.3 & \textbf{24.5} \\ 
\bottomrule
\end{tabular}
\end{adjustbox}
\caption{\label{tab:multiple-e} \ours{} 15B attains high competency in coding performance across a broad range of programming languages. Results for Mistral are from our runs of Mistral in the same setting as \ours{}. }
\end{table}

\subsection{Multilingual}

We demonstrate the outstanding multilingual ability of \ours{} 15B using four widely-studied benchmarks in previous works that cover a diverse range of high to low resource natural languages. For classification we use accuracy as the metric; for generative tasks, we use exact match; and for machine translation, we evaluate using the \texttt{sacreBLEU} \citep{post2018} implementation of \texttt{BLEU}~\citep{papineni2002}, using \texttt{spm-flores-101} tokenization to obtain spBLEU scores. 

\textbf{1. Classification:}  Cross-lingual Choice of Plausible Alternatives (XCOPA) \citep{xcopa} tests causal commonsense reasoning in 11 languages 

We compare \ours{} 15B to existing multilingual language models: XGLM \citep{lin2022fewshot} , mGPT \citep{shliazhko2022mgpt}, and BLOOM \citep{workshop2023bloom}. 
XGLM and mGPT are models specially trained to have improved multilingual ability by up-sampling the presence of non-English languages in the training data. 
In contrast, BLOOM, like \ours{}, is a general purpose language model that was trained on a combination of English, multilingual, and code data. 
In Table \ref{tab:xcopa}, we clearly see that \ours{} achieves the best performance amongst all models -- realizing almost a 12\% improvement in the four-shot setting. 

\begin{table}[H]
\centering
\begin{adjustbox}{max width=\textwidth, scale=1.0}
\begin{tabular}{lcrcccccccccccccc}
\toprule
Mode & Model & Size & ET & HT & ID & IT  & QU & SW & TA  & TH & TR & VI & ZH & AVG \\
\midrule
 \multirow{4}{*}{Zero-Shot} & BLOOM &   176B & - & - &  $57.5^*$ &	- &	 -  &	$59.5^*$ &	 $54.7^*$ &	- &	- &  $58.2^*$ &  $57.7^*$ & -     \\
 & XGLM & 7.5B & 57.6	& 57.0	& 59.0 &	49.2 &	52.4&	55.0 &	55.6	& 57.8& 	55.0 &	59.0 &	53.6 &	55.6 \\
 & mGPT &   13B & 49.8 &	50.4	& 63.4	& 61.6	& 50.4	& 57.6 &	57.0 & 54.0 &	58.2	& 60.4	& 54.6 & 56.1  \\
  & \ours{} & 15B & 62.8 & 47.4 & 66.6 & 67.0 & 53.8 & 50.4 & 62.0 & 59.6 & 57.4 & 65.2 & 62.2 & \textbf{59.5}   \\
  \midrule
\multirow{3}{*}{4-Shot} & XGLM &   7.5B & 64.7	&60.4	&67.3	& 64.0&	50.0&	61.8&	56.7&	61.5&	60.1 &	68.5 & 59.9	& 61.4   \\
 & mGPT & 13B & 48.6&	48.6	&62.6&	60.8&	50.6&	56.6&	55.4&	54.8&	57.4	&61.8&	58.4&56.0            \\
 & \ours{} & 15B & 72.9 & 52.8 & 79.6 & 79.2 & 50.2 & 52.2 & 72.8 & 66.6 & 77.2 & 78.6 & 76.0 & \textbf{68.9}  \\
\bottomrule
\end{tabular}
\end{adjustbox}
\caption{\label{tab:xcopa} Comparison of \ours{} 15B against existing large language models on XCOPA under the zero- and four-shot setting. Our reported results for XGLM are from the runs of the model in \citep{shliazhko2022mgpt} given that we use the same prompt template used by mGPT. The values marked with $*$ are read from figures in \citep{workshop2023bloom}. 
}
\end{table}

\textbf{2. Generation:} We consider two generative tasks: TyDiQA-GoldP~\citep{tydiqa} and Multilingual Grade School Math (MGSM)~\citep{shi2022language}. TyDiQA-GoldP is a question answering task while MGSM evaluates the arithmetic reasoning ability of language models in 10 languages.

In comparing the performance of \ours{} 15B on TyDiQA-GoldP to a range of models, Table \ref{tab:tydiqa_internal} shows that \ours{} 15B achieves the best performance. Impressively, \ours{} 15B is able to significantly improve upon the next best model, PaLM 62B-cont. 

\begin{table}[H]
\centering
\begin{adjustbox}{max width=\textwidth, scale=1.0}
\begin{tabular}{crrrrrrrrrrrr}
\toprule
Model & Size & AR & BN & FI & ID & KO & RU & SW & TE & AVG  \\
\midrule
\multirow{2}{*}{PaLM}& 62B & 31.2 & 42.5 & 41.7 & 41.6	& 49.3 & 29.2 &	58.1 & 30.6 &	40.5    \\
 & 62B-cont & 39.4	& 48.7 & 44.0 &	49.2	&52.5&	35.6 &	60.9&	35.3	& 45.7    \\
 LLaMA-2 & 13B & -& -& -& -& -& -& -& -& 33.2 \\
 Baichuan-2 & 13B & -& -& -& -& -& -& -& -& 30.8\\
 QWEN & 14B & -& -& -& -& -& -& -& -& 39.8 \\
\ours{} & 15B & 39.1 & 55.8 &52.2 & 54.5 & 55.1 & 37.8 & 54.5 & 55.0 & \textbf{50.5} \\
\bottomrule
\end{tabular}
\end{adjustbox}
\caption{\label{tab:tydiqa_internal} Comparative results in the one-shot setting on TyDiQA-GoldP. Results for LLaMA-2 13B, Baichuan-2 13B and QWEN 14B are taken from \citep{chen2024orion14b}.  }
\end{table}

Further demonstrating the impressive multilingual ability of \ours{} 15B, Table \ref{tab:mgsm_all} shows the performance on MGSM. 
We report using the English chain-of-thought setting introduced in \citep{shi2022language} where all chain of thought explanations are presented to the model in English rather than in the language of the task. On this challenging task which assesses the intersection of mathematical and multilingual ability, \ours{} 15B achieves the best performance amongst compared models and improves upon the closest score by nearly 30\%.

\begin{table}[H]
\centering
\begin{adjustbox}{max width=\textwidth, scale=1.0}
\begin{tabular}{lcrrrrrrrrrrrrrrr}
\toprule
Mode & Model & Size & DE & FR & ES & RU & ZH & JA & TH & TE & BN & SW & AVG  \\
\midrule
 \multirow{1}{*}{Native-COT}  & \multirow{1}{*}{PaLM} 
 &   62B 	&24.0  &	24.0 &	26.0 &	22.8 &	24.8	& 14.8 & 18.0 & 11.6 &	13.6 & 9.6&	18.9	\\
\midrule
\multirow{3}{*}{English-COT} & \multirow{1}{*}{PALM} & 62B-cont & 44.8 & 39.2 & 44.4 & 36.8 & 33.6 & 24.0 & 28.0 & 19.6 & 28.0 & 21.2 & 32.0 \\
& \multirow{1}{*}{Mistral} & 7B & 33.2 & 35.2 & 35.6 & 35.2 & 33.2 & 18.8 & 10.0 & 0.0 & 8.0 & 9.2 & 21.8 \\
& \multirow{1}{*}{\ours{}} & 15B & 46.8 & 46.0 & 50.0&45.6&40.0&40.0&43.6&41.6&43.6&16.0 & \textbf{41.3} \\
\bottomrule
\end{tabular}
\end{adjustbox}
\caption{\label{tab:mgsm_all} Eight-shot accuracy results on MGSM. Results for Mistral are from our runs of Mistral in the same setting as \ours{}. }
\end{table}

\textbf{3. Machine Translation:} We additionally evaluate the translation ability of our models through the FLORES-101 \citep{goyal2021} benchmark. The ability to translate between languages is a good test of the model's ability to relate and understand semantic relationships between languages.

As seen in Table~\ref{tab:flores_ext}, \ours{} 15B heftily outperforms both LLaMA-2 13B and Baichuan-2 13B -- improving upon their performance by 90.2\% and 44.1\% respectively. \ours{} 15B does not solely perform well on translating from Chinese into English but is able to attain impressive results on the direct translation of Chinese into other languages. This ability highlights the strong understanding that \ours{} 15B has across a broad spectrum of natural languages.

\begin{table}[H]
\centering
\begin{adjustbox}{max width=\textwidth, scale=1.0}
\begin{tabular}{lcrrrrrrrr}
\toprule
  & Size & ZH-EN & ZH-FR & ZH-ES & ZH-AR & ZH-RU & ZH-JA & ZH-DE & AVG  \\
\midrule
\multirow{1}{*}{LLaMA-2} 
& 13B & 25.4 &19.2 &17.5& 1.4& 10.3& 0.1 &11.1 &12.2  \\
\midrule
\multirow{1}{*}{Baichuan-2} 
& 13B & 30.6 & 22.1 & 17.3 & 2.4 & 14.2 & 11.6 & 14.5 & 16.1 \\
\midrule
\ours{} & 15B & 34.0 & 28.1 & 21.3 & 16.8 & 21.2 & 23.1 &	18.1 & \textbf{23.2} \\ 
\bottomrule
\end{tabular}
\end{adjustbox}
\caption{\label{tab:flores_ext} Eight-shot results on Flores sub-tasks translating out of Chinese. All results for external models were obtained from \citep{yang2023baichuan}}
\end{table}

\section{Conclusion}

We present \ours{} 15B, a decoder-only transformer-based large language model.
It is trained on 8 trillion tokens spanning English, 53 additional natural languages as well as 43 programming languages. \ours{} 15B exhibits the strongest multilingual performance of any general purpose language model at its scale -- even outperforming models specialized for the multilingual domain.
\ours{} demonstrates that pre-training sets for large language models can continue to be scaled up even further in order to improve the abilities of models. 

\bibliographystyle{plainnat}
\bibliography{references}

\FloatBarrier

\appendix

\section*{Supplementary Materials} \label{supplementary_material}

\begin{table}[H]
\centering
\begin{tabular}{lrccccc} 
\toprule
    & Size & Humanities & Social sciences & STEM & Other & Average       \\ 
\midrule
\ours{} & 15B & 69.2 & 74.1 & 53.4 & 67.5 & 64.2 \\ 
\bottomrule
\end{tabular}
\caption{\label{tab:mmlu} Per-category breakdown accuracy for MMLU }
\end{table}

\end{document}
